\Needspace{8\baselineskip}\originalchapter{2006 作业四: 横向散射与空间电荷波}
题源: 5 Oct 2006, due 17 Oct.
\begin{exercise}\label{ex:06-4-1}
原题1. 快离子束能量E、速度v0远大于等离子体电子离子热速, ne=ni=n0. (i)求横向扩散频率, 题印 $\dd\langle v_\perp^2\rangle/\dd t=-\nu_\perp v_0^2$, 用E(eV)、n、ln\ensuremath{\Lambda}、相对质子质量电荷表示, 忽略1/ln\ensuremath{\Lambda}修正; (ii)与课程资源链接NRL Formulary pp31--32比较.
\end{exercise}
\begin{solution}
扩散方差增长, 原定义的负号有误; 本解取 $\dd\langle v_\perp^2\rangle/\dd t=+\nu_\perp v_0^2$、\ensuremath{\nu}\ensuremath{\perp}>0. 对靶s, $K_s=q_bq_s/(4\pi\epsilon_0)$, 小角冲量 $\Delta v_\perp=2K_s/(m_bv_0b)$. 乘 $n_sv_0\,2\pi b\dd b$, 积两截止得
\begin{equation}
 \nu_\perp=\sum_s\frac{8\pi n_sK_s^2\ln\Lambda_s}{m_b^2v_0^3}
 =1.80325\times10^{-13}\frac{Z_b^2\sum_s n_sZ_s^2\ln\Lambda_s}{\sqrt{A_b}\,E_{\rm eV}^{3/2}}\ \mathrm{s}^{-1},
\end{equation}
n用m$^{-3}$、mb=Abmp. 对源题ne=ni=n0和单电荷离子, 和为2n0ln\ensuremath{\Lambda}. 若离子Zi\ensuremath{\ne}1, 准中性应ne=Zini, 不能同时维持原ne=ni.

(ii)对照时须以NRL的deflection频率及其热速、能量与密度单位匹配: 改n为cm$^{-3}$后前因子变 $1.80325\times10^{-7}$, 不能直接与减速频率或单横分量扩散频率比较, 后者少因子2. 课程实际链接2002版DTIC文件本次访问返回代理403, 未取得指定页; 因而不声称已完成逐式NRL原页核对. 上式为独立冲量推导, 此外部页核对在覆盖清单列为未取得参考, 不编造其文字或数值表.
\end{solution}
\begin{exercise}\label{ex:06-4-2}
原题2. 一维冷电子束与不动中性离子背景, 连续性、$mnD\mathbf v/Dt=-en\mathbf E$, Gauss $\epsilon_0\partial_xE=-e(n_e-n_i)$. ne=n0+n1,ni=n0,ve=Ve0+v1,E=E1. (a)完整线性方程; (b)以 $\exp(\ii\omega t-\ii kx)$ 求并画k(\ensuremath{\omega}), 解释快慢波; (c)求每支位移加束电流 $J_x=\epsilon_0\partial_tE_1+J_b$, 为何可忽略磁场?
\end{exercise}
\begin{solution}
(a) $(\partial_t+V_0\partial_x)n_1+n_0\partial_xv_1=0$, $m(\partial_t+V_0\partial_x)v_1=-eE_1$, $\epsilon_0\partial_xE_1=-en_1$. (b)联立 $(\omega-kV_0)^2=\omega_p^2$, 故 $k=(\omega\mp\omega_p)/V_0$. \ensuremath{\omega}>\ensuremath{\omega}p的同向传播时小k支相速大于V0称快, 大k支小于V0称慢.
\begin{center}\begin{tikzpicture}\begin{axis}[width=.6\textwidth,height=4cm,xlabel={$\omega/\omega_p$},ylabel={$kV_0/\omega_p$},xmin=0,xmax=3,ymin=-1,ymax=4,legend style={font=\small}]
\addplot[thick,structurecolor,domain=0:3]{x-1};\addlegendentry{快支}\addplot[thick,dashed,domain=0:3]{x+1};\addlegendentry{慢支}\end{axis}\end{tikzpicture}\end{center}
(c)扰动束电流 $J_b=-e(n_0v_1+V_0n_1)$. 连续性给 $J_b=(\omega/k)(-en_1)=-\ii\omega\epsilon_0E_1$在题时间约定下, 正好抵消位移项 $+\ii\omega\epsilon_0E_1$. 所以纵波总Amp\`ere源零, 可取B1=0; 不是因电子电流本身为零.
\end{solution}
\begin{exercise}\label{ex:06-4-3}
原题3. 冷均匀等离子体电子波频率与k无关, 加压力后研究色散. (a)对连续性、$m_enD\mathbf v/Dt=-en\mathbf E-\nabla p$, $pn^{-\gamma}=p_0n_0^{-\gamma}$、Gauss在线性静背景展开; (b) $\exp(-\ii\omega t+\ii kx)$ 求非平凡解条件.
\end{exercise}
\begin{solution}
$-\ii\omega n_1+\ii kn_0v_1=0$, $-\ii\omega m_en_0v_1=-en_0E_1-\ii kp_1$, $p_1=\gamma p_0n_1/n_0$, $\ii k\epsilon_0E_1=-en_1$. 消元 $\omega^2=\omega_p^2+\gamma p_0k^2/(m_en_0)$. 对高相速一维热运动取\ensuremath{\gamma}3得Bohm--Gross; 若指定等温\ensuremath{\gamma}1是另一闭合. 无压力p0=0才恢复与k无关的冷频率.
\end{solution}
